Learned construction policies are proposed for metro network expansion, but they are rarely compared with evolutionary multi-objective search at an equal simulation budget, or scored on every objective. We compare a weight-conditioned graph reinforcement learning policy with NSGA-II, scalarised genetic algorithms, tabular Q-learning, a greedy rule and random search on the public Xi'an and Amsterdam benchmark grids, with about $10^4$ simulated lines per method and ten seeds. We replace the benchmark's summed percentage reward by the true share of demand served, score fairness by a generalised Gini welfare that is zero when nobody is served, and measure accessibility as the trip-end mass newly reached. On Xi'an, the search archive of the learned policy is comparable with that of NSGA-II (hypervolume \XianRLArchHv\ against \XianNSGAArchHv, $p=\XianRLvsNSGAArchHvPHolm$), but its 15 deployable lines are worse (expected utility \XianRLEum\ against \XianNSGAEum, $p=\XianRLvsNSGAPolicyEumPHolm$). On the larger Amsterdam grid NSGA-II is better on both measures. The policy produces nearly the same line for different weights. In a leave-one-line-out test on Xi'an, NSGA-II recovers the real Line~2 about twice as well as a random-search baseline, and the learned policy does not differ from random. A warm start from another city helped in one direction (Amsterdam to Xi'an, \TrAmsXianFullWins\ of 5 seeds) and not in the other, and five pairs are too few to establish either effect. For this benchmark, NSGA-II is the stronger choice at equal budget; the learned policy is not yet a replacement.
